\begin{lemma}[分块预解式的 Schur 补]
设 $H$ 为有界算符，$P,Q=I-P$ 为互补正交投影。若 $z-QHQ$ 与 $z-H$ 可逆，则 $P(z-H)^{-1}P=[z-PHP-PHQ(z-QHQ)^{-1}QHP]^{-1}$，其中右端作用在 $P\Hil$ 上。有限维矩阵是此结论的特例。无界算符需另外保证所用分块及乘积的共同定义域。
\end{lemma}
\begin{proof}
要计算目标空间上的预解式，先在\eqnrefs{eq:feshbach-block-matrix-dp68} 的分块方程中消去补空间分量，再求剩余目标块的逆。这里使用补块的可逆性，所得 Schur 补关系适用于任意满足上述假设的耦合强度。

记
\begin{equation*}
 z-H=
 \begin{pmatrix}A&B\\C&D\end{pmatrix},
 \quad
 A=z-PHP,
 \quad B=-PHQ,
 \quad C=-QHP,
 \quad D=z-QHQ.
\end{equation*}
设其逆矩阵为
\begin{equation*}
 \begin{pmatrix}X&Y\\Z&W\end{pmatrix}.
\end{equation*}
由下左块方程
\begin{equation*}
 CX+DZ=0
\end{equation*}
得到
\begin{equation*}
 Z=-D^{-1}CX.
\end{equation*}
代入上左块 $AX+BZ=I_P$，
\begin{align*}
 (A-BD^{-1}C)X&=I_P,\\
 X&=(A-BD^{-1}C)^{-1}.
\end{align*}
而 $X=P(z-H)^{-1}P$。恢复四个分块以后即得正文\eqnrefs{eq:feshbach-schur-proof-r10}。
\end{proof}

\begin{proposition}[Schrieffer--Wolff 二阶矩阵元]
在有限维模型中设 $H=H_{\rm d}+H_{\rm od}$，目标块与补块的未扰动谱相隔正谱隙，$S$ 取反 Hermitian 的非对角解 $[S,H_{\rm d}]=-H_{\rm od}$。则目标块的二阶修正为 $\frac12\sum_{q\in Q}H_{pq}H_{qp^\prime}[(E_p-E_q)^{-1}+(E_{p^\prime}-E_q)^{-1}]$，余项按弱混合参数取三阶。
\end{proposition}
\begin{proof}
Baker--Campbell--Hausdorff 展开给出

\begin{align*}
 e^SHe^{-S}
 &=H+[S,H]+\frac12[S,[S,H]]+\order(S^3H)\\
 &=H_{\rm d}+H_{\rm od}+[S,H_{\rm d}]+[S,H_{\rm od}]
 +\frac12[S,[S,H_{\rm d}]]+\order(S^3H).
\end{align*}
由正文\eqnrefs{eq:sw-sylvester-r15}，$[S,H_{\rm d}]=-H_{\rm od}$，因此
\begin{align*}
 e^SHe^{-S}
 &=H_{\rm d}+[S,H_{\rm od}]
 -\frac12[S,H_{\rm od}]+\order(S^3H)\\
 &=H_{\rm d}+\frac12[S,H_{\rm od}]+\order(S^3H),
\end{align*}
即正文\eqnrefs{eq:sw-bch-second-r15}。

令
\begin{equation*}
 H_{\rm d}|p\rangle=E_p|p\rangle,
 \qquad
 H_{\rm d}|q\rangle=E_q|q\rangle,
\end{equation*}
其中 $p,p'\in P$、$q\in Q$。在 $\langle q|\cdots|p\rangle$ 上取正文\eqnrefs{eq:sw-sylvester-r15} 的矩阵元，
\begin{align*}
 0&=H_{qp}+(E_p-E_q)S_{qp},\\
 S_{qp}&=\frac{H_{qp}}{E_q-E_p},
 \qquad
 S_{pq}=\frac{H_{pq}}{E_p-E_q}.
\end{align*}
于是
\begin{align*}
 \langle p|[S,H_{\rm od}]|p'\rangle
 &=\sum_{q\in Q}
 \left(S_{pq}H_{qp'}-H_{pq}S_{qp'}\right)\\
 &=\sum_{q\in Q}H_{pq}H_{qp'}
 \left[
 \frac1{E_p-E_q}
 +\frac1{E_{p'}-E_q}
 \right].
\end{align*}
乘上正文\eqnrefs{eq:sw-bch-second-r15} 中的 $1/2$，并加回 $P H_{\rm d}P$，得到正文\eqnrefs{eq:sw_arbitrary_target_subspace_matrix_elements}。
\end{proof}

\begin{theorem}[自伴预解式的精确范数]
若 $A$ 为自伴算符且 $z\notin\sigma(A)$，则 $\|(z-A)^{-1}\|=\operatorname{dist}(z,\sigma(A))^{-1}$。该结论同时适用于离散谱与连续谱。
\end{theorem}
\begin{proof}
正文\eqnrefs{eq:resolvent-spectral-resolution-dp68} 把自伴算符的预解式写成谱测度积分。对任意归一化态取范数平方，可以把算符范数问题化成正概率测度下的标量函数 $|z-\lambda|^{-2}$ 的上确界，从而得到正文\eqnrefs{eq:resolvent-norm-spectral-dp21}。

对归一化 $|\psi\rangle$ 定义
\begin{equation*}
 \dd\mu_\psi(\lambda)
 =\langle\psi|\dd E_A(\lambda)|\psi\rangle,
 \qquad
 \int\dd\mu_\psi=1.
\end{equation*}
于是
\begin{align*}
 \|(z-A)^{-1}|\psi\rangle\|^2
 &=\int_{\sigma(A)}
 \frac{\dd\mu_\psi(\lambda)}{|z-\lambda|^2}\\
 &\le
 \sup_{\lambda\in\sigma(A)}\frac1{|z-\lambda|^2}.
\end{align*}
对所有归一化态取上确界，得到
\begin{equation*}
 \|(z-A)^{-1}\|
 \le\frac1{\operatorname{dist}[z,\sigma(A)]}.
\end{equation*}
为得到反向不等式，记 $d=\operatorname{dist}(z,\sigma(A))$。任给 $\epsilon>0$，可选谱点 $\lambda_0$ 使 $|z-\lambda_0|<d+\epsilon/2$，再选其足够小的邻域 $U$，使 $|z-\lambda|<d+\epsilon$ 对 $\lambda\in U$ 成立。由于 $\lambda_0$ 在谱测度的支撑上，$E_A(U)\ne0$。从其像中取归一化态，便有
\[
 \|(z-A)^{-1}\psi\|^2
 =\int_U\frac{\dd\mu_\psi(\lambda)}{|z-\lambda|^2}
 \ge\frac1{(d+\epsilon)^2}.
\]
令 $\epsilon\downarrow0$ 即得反向界，与上界合并后得到等号。这一论证不要求存在最近本征态，因而也覆盖连续谱。令 $A=QHQ$ 即得\eqnrefs{eq:resolvent-norm-spectral-dp21}。
\end{proof}

\begin{proposition}[有谱隙的静态有效算符展开]
设补块 $QHQ$ 自伴，$E_*\in\RR$ 与补块谱的距离为 $\Delta_Q>0$，$|z-E_*|<\Delta_Q$，且分块耦合有界。则补块预解式可在 $E_*$ 处作算符范数收敛的 Neumann 展开；静态首项为
\[
H_{\rm eff}^{(0)}=PHP+PHQ(E_*-QHQ)^{-1}QHP.
\]
精确有效算符为 $H_{\rm eff}(z)=H_{\rm eff}^{(0)}+R(z)$，其中
\[
\|R(z)\|\le
\frac{\|PHQ\|\,\|QHP\|\,|z-E_*|}
{\Delta_Q(\Delta_Q-|z-E_*|)}.
\]
当 $PHP=E_*P$ 或目标块自身的谱展宽为 $O(\epsilon_{\rm mix}^2\Delta_Q)$ 时，靠近该目标块的极点才给出下述四阶冻结误差。
\end{proposition}
\begin{proof}
精确 Feshbach--Schur 有效算符为
\[
H_{\rm eff}(z)
=
PHP+\Sigma_P(z),
\qquad
\Sigma_P(z)
=
PHQ\,(z-QHQ)^{-1}QHP.
\]
这里 \(P\) 可以是任意有限维目标子空间，并不要求只保留一个本征态。设目标能窗以参考能量 \(E_*\) 为中心，并与 \(QHQ\) 的谱保持最小距离
\[
\Delta_Q
\equiv
\operatorname{dist}\!\left(E_*,\sigma(QHQ)\right)>0.
\]
若实际极点满足 \(|z-E_*|<\Delta_Q\)，则
\[
z-QHQ
=
(E_*-QHQ)
\left[
I+(z-E_*)(E_*-QHQ)^{-1}
\right].
\]
定义
\[
X\equiv
(z-E_*)(E_*-QHQ)^{-1}.
\]
由预解式范数
\[
\|(E_*-QHQ)^{-1}\|
\le\frac1{\Delta_Q}
\]
可知
\[
\|X\|
\le\frac{|z-E_*|}{\Delta_Q}<1,
\]
所以在目标能窗内可以作算符 Neumann 展开：
\begin{align*}
(z-QHQ)^{-1}
&=
\left(I+X\right)^{-1}(E_*-QHQ)^{-1}\\
&=
\left(I-X+X^2-\cdots\right)
(E_*-QHQ)^{-1}\\
&=
(E_*-QHQ)^{-1}
-(z-E_*)(E_*-QHQ)^{-2}
+\mathcal O\!\left(
\frac{|z-E_*|^2}{\Delta_Q^3}
\right).
\end{align*}
因此
\begin{align*}
\Sigma_P(z)
&=
PHQ(E_*-QHQ)^{-1}QHP\\
&\quad
-(z-E_*)
PHQ(E_*-QHQ)^{-2}QHP
+\cdots .
\end{align*}

若 \(V\equiv PHQ+QHP\) 是弱混合，并定义
\[
\epsilon_{\rm mix}
\equiv
\frac{\|PHQ\|}{\Delta_Q}\ll1,
\]
则第一项满足
\[
\|\Sigma_P(E_*)\|
\le
\frac{\|PHQ\|^2}{\Delta_Q}
=
\epsilon_{\rm mix}^2\Delta_Q.
\]
一般目标块不能仅由弱混合推出 $|z-E_*|$ 为二阶，因为 $PHP$ 可以自带有限谱展宽。预解恒等式首先给出对任意允许 $z$ 都成立的估计
\[
\|\Sigma_P(z)-\Sigma_P(E_*)\|
\le\frac{\|PHQ\|\,\|QHP\|\,|z-E_*|}
{\Delta_Q(\Delta_Q-|z-E_*|)}.
\]
这里用到了自伴补块的预解式范数界。再假设全算符自伴且 $PHP=E_*P$，或更一般地
\[
\|PHP-E_*P\|=O(\epsilon_{\rm mix}^2\Delta_Q).
\]
若实极点 $z$ 位于 $|z-E_*|\le\Delta_Q/2$ 的目标能窗，取归一化有效本征向量便得
\[
|z-E_*|\le\|PHP-E_*P\|+\|\Sigma_P(z)\|
\le\|PHP-E_*P\|+2\epsilon_{\rm mix}^2\Delta_Q.
\]
所以在这一附加条件下，精确极点相对参考能量的移动是二阶，
\[
|z-E_*|
=
\mathcal O(\epsilon_{\rm mix}^2\Delta_Q).
\]
代回第二项可得
\[
\left\|
(z-E_*)
PHQ(E_*-QHQ)^{-2}QHP
\right\|
=
\mathcal O(\epsilon_{\rm mix}^4\Delta_Q),
\]
所以在二阶精度下，自能中的能量依赖可以冻结在 \(E_*\)：
\[
H_{\rm eff}^{(2)}
=
PHP
+
PHQ(E_*-QHQ)^{-1}QHP
+
\mathcal O(\epsilon_{\rm mix}^4\Delta_Q).
\]

正文\eqnrefs{eq:feshbach-second-order-static-dp68} 是一维目标子空间的特例。取
\[
P=|p\rangle\langle p|,
\qquad
PHP=E_pP,
\qquad
E_*=E_p,
\]
并对 \(QHQ\) 插入本征态完备关系
\[
Q
=
\sum_{q\in Q}|q\rangle\langle q|,
\qquad
QHQ|q\rangle=E_q|q\rangle,
\]
得到
\begin{align*}
\langle p|\Sigma_P(E_p)|p\rangle
&=
\sum_{q\in Q}
\frac{
\langle p|H|q\rangle
\langle q|H|p\rangle
}{E_p-E_q}\\
&=
\sum_{q\in Q}
\frac{|H_{pq}|^2}{E_p-E_q}.
\end{align*}
于是
\[
z
=
E_p
+
\sum_{q\in Q}
\frac{|H_{pq}|^2}{E_p-E_q}
+
\mathcal O(\epsilon_{\rm mix}^4\Delta_Q),
\]
即正文\eqnrefs{eq:feshbach-second-order-static-dp68}。因此“二阶能量修正”不是从非简并微扰论类比得到，而是精确 Schur 补在谱隙与弱混合控制下的静态展开；若目标子空间简并或近简并，应保留整个 \(P\) 块而不能先选单个 \(|p\rangle\)。
\end{proof}

\begin{proposition}[盒归一化的连续谱极限]
周期立方盒的动量为 $\mathbf p_{\mathbf n}=2\pi\hbar\mathbf n/L$。对足够快衰减的连续函数 $f$，$L^{-3}\sum_{\mathbf n}f(\mathbf p_{\mathbf n})\to\int f(\mathbf p)\,\dd^3p/(2\pi\hbar)^3$；相应完备性核在缓增分布意义下趋于 $\delta^{(3)}$。
\end{proposition}
\begin{proof}
对随动量变化足够平滑的测试函数 $F(\mathbf p)$，

\begin{align*}
 \sum_{\mathbf n}F(\mathbf p_{\mathbf n})
 &=\sum_{\mathbf n}
 \left(\frac{2\pi\hbar}{L}\right)^3
 \frac{F(\mathbf p_{\mathbf n})}{(2\pi\hbar/L)^3}\\
 &\longrightarrow
 \frac{L^3}{(2\pi\hbar)^3}
 \int\dd^3 p\,F(\mathbf p).
\end{align*}
这给出第一条极限。

盒归一化平面波为
\begin{equation*}
 \langle\mathbf r|\mathbf n\rangle
 =L^{-3/2}\ee^{\ii\mathbf p_{\mathbf n}\cdot\mathbf r/\hbar},
\end{equation*}
离散完备性
$\sum_{\mathbf n}|\mathbf n\rangle\langle\mathbf n|=I$
在位置表象中变成
\begin{equation*}
 \frac1{L^3}\sum_{\mathbf n}
 \ee^{\ii\mathbf p_{\mathbf n}\cdot(\mathbf r-\mathbf r')/\hbar}
 =\langle\mathbf r|I|\mathbf r'\rangle.
\end{equation*}
将刚得到的求和--积分极限用于左端，
\begin{equation*}
 \frac1{L^3}\sum_{\mathbf n}
 \ee^{\ii\mathbf p_{\mathbf n}\cdot\Delta\mathbf r/\hbar}
 \longrightarrow
 \int\frac{\dd^3 p}{(2\pi\hbar)^3}
 \ee^{\ii\mathbf p\cdot\Delta\mathbf r/\hbar}
 =\delta^{(3)}(\Delta\mathbf r),
\end{equation*}
即正文\eqnrefs{eq:box-to-continuum-dp21} 的第二条关系。
\end{proof}

\begin{lemma}[Grassmann 高斯行列式]
采用本章固定的 Berezin 测度顺序，有限个独立 Grassmann 变量的高斯积分满足 $\mathcal Z_G[0,0]=\int\prod_i\dd\bar\theta_i\dd\theta_i\,\exp(-\bar\theta M\theta)=\det M$。此式不要求 $M$ 可逆。
\end{lemma}
\begin{proof}
Grassmann 高斯积分产生行列式的原因可以直接从“积分抽取最高次系数”看出。因为每个 $\theta_i$ 和 $\bar\theta_i$ 都平方为零，指数
\begin{equation*}
 \ee^{-\bar\theta M\theta}
 =\sum_{r=0}^{n}\frac{(-1)^r}{r!}(\bar\theta M\theta)^r
\end{equation*}
是有限多项式。测度含全部 $2n$ 个变量，所以积分后只有 $r=n$ 项可能留下。

展开 $(\bar\theta M\theta)^n$ 时，只要两个行指标或两个列指标重复，相应单项式就含某个 Grassmann 变量的平方而消失。因此非零项恰由 $1,\ldots,n$ 的排列标记。把变量重排回正文固定的测度次序会产生排列符号，而第一组指标的 $n!$ 种排列抵消展开中的 $1/n!$。于是
\begin{align*}
 \mathcal Z_G[0,0]
 &=\sum_{\sigma\in S_n}\operatorname{sgn}(\sigma)
 \prod_{i=1}^{n}M_{i\sigma(i)}\\
 &=\det M,
\end{align*}
最后一行正是行列式的 Leibniz 定义，从而得到正文\eqnrefs{eq:grassmann-gaussian-determinant-dp68}。
\end{proof}

\begin{lemma}[含 Grassmann 源的高斯积分]
若矩阵 $M$ 可逆，且源与测度次序如本章定义，则 $\mathcal Z_G[\xi,\bar\xi]=\det M\exp(\bar\xi M^{-1}\xi)$。
\end{lemma}
\begin{proof}
加入源以后，Grassmann 高斯积分与普通高斯一样可以完成平方，只是变量和源的次序必须保持固定。令
\begin{equation*}
 \theta'=\theta-M^{-1}\xi,
 \qquad
 \bar\theta'=\bar\theta-\bar\xi M^{-1},
\end{equation*}
也就是 $\theta=\theta'+M^{-1}\xi$、$\bar\theta=\bar\theta'+\bar\xi M^{-1}$。代回指数并保持因子次序，有
\begin{align*}
 &-\bar\theta M\theta+\bar\xi\theta+\bar\theta\xi\\
 &=-\bar\theta'M\theta'
 -\bar\theta'\xi-\bar\xi\theta'-\bar\xi M^{-1}\xi
 +\bar\xi\theta'+\bar\xi M^{-1}\xi
 +\bar\theta'\xi+\bar\xi M^{-1}\xi\\
 &=-\bar\theta'M\theta'+\bar\xi M^{-1}\xi.
\end{align*}
Berezin 积分对这种常数 Grassmann 平移保持不变，因此
\begin{align*}
 \mathcal Z_G[\xi,\bar\xi]
 &=\ee^{\bar\xi M^{-1}\xi}
 \int\prod_i\dd\bar\theta_i'\dd\theta_i'\,
 \ee^{-\bar\theta'M\theta'}\\
 &=\mathcal Z_G[0,0]\ee^{\bar\xi M^{-1}\xi},
\end{align*}
即正文\eqnrefs{eq:grassmann-source-gaussian-dp68}。
\end{proof}

\begin{corollary}[Grassmann 二点收缩]
在同一左右导数及源顺序约定下，归一化高斯生成函数的二阶源导数在零源处给 $(M^{-1})_{ij}$，具体次序为\eqnrefs{eq:grassmann-inverse-kernel-dp68}。
\end{corollary}
\begin{proof}
正文已经固定了源的次序，因此二点收缩只需按同一约定求两次 Grassmann 导数。先对 $\bar\xi_i$ 取左导数，
\begin{equation*}
 \frac{\overrightarrow\partial\mathcal Z_G}{\partial\bar\xi_i}
 =\mathcal Z_G(M^{-1})_{ik}\xi_k.
\end{equation*}
再从右边对 $\xi_j$ 求导，并令源为零，指数因子回到 $1$，从而
\begin{equation*}
 \left.
 \frac{1}{\mathcal Z_G[0,0]}
 \frac{\overrightarrow\partial\mathcal Z_G}{\partial\bar\xi_i}
 \frac{\overleftarrow\partial}{\partial\xi_j}
 \right|_{0}
 =(M^{-1})_{ij}.
\end{equation*}
这就是正文\eqnrefs{eq:grassmann-inverse-kernel-dp68}。若交换源或导数的次序，会因 Grassmann 奇偶性出现额外负号，所以以后把这种带 Grassmann 源的高斯积分用于费米场时，必须同时固定源、导数和测度的次序。
\end{proof}

\begin{lemma}[Berezin 测度的逆行列式律]
对可逆线性变换 $\theta^\prime=A\theta$，有 $\dd^n\theta^\prime=(\det A)^{-1}\dd^n\theta$，其中两边使用同一变量排列顺序。
\end{lemma}
\begin{proof}
普通变量的 Jacobian 来自体积元，而 Berezin 测度由最高次系数的归一化决定，所以变换方向恰好相反。先看一维 $\theta'=a\theta$。若写 $\dd\theta'=C\dd\theta$，则\eqnrefs{eq:berezin-definition-dp68} 要求
\begin{equation*}
 1=\int\dd\theta'\,\theta'
 =Ca\int\dd\theta\,\theta=Ca,
\end{equation*}
故 $C=a^{-1}$。

多维时由\eqnrefs{eq:berezin-linear-transform-dp68} 有
\begin{equation*}
 \theta_1'\cdots\theta_n'
 =A_{1i_1}\cdots A_{ni_n}\theta_{i_1}\cdots\theta_{i_n}.
\end{equation*}
只有 $(i_1,\ldots,i_n)$ 为一个排列时该项非零；重排到固定次序所产生的符号正好组成行列式，因此
\begin{equation*}
 \theta_1'\cdots\theta_n'=(\det A)\theta_1\cdots\theta_n.
\end{equation*}
若 $\dd^n\theta'=C\dd^n\theta$，最高次系数的 Berezin 归一化给出 $C\det A=1$，于是得到正文\eqnrefs{eq:berezin-jacobian-dp21}。
\end{proof}

\begin{proposition}[外微分的幂零性与 Leibniz 规则]
光滑 $p$-形式 $\omega$ 与任意形式 $\eta$ 满足 $\dd^2\omega=0$ 和 $\dd(\omega\wedge\eta)=\dd\omega\wedge\eta+(-1)^p\omega\wedge\dd\eta$。
\end{proposition}
\begin{proof}
外微分的两个基本性质都来自同一件事：普通偏导对光滑函数可以交换，而基一形式的楔积在交换次序时变号。对 $p$ 阶形式 $\omega$ 再作用一次外微分，
\begin{equation*}
 \dd^2 \omega
 =\frac1{p!}\partial_\rho\partial_\nu\omega_{\mu_1\cdots\mu_p}
 \dd x^\rho\wedge\dd x^\nu\wedge
 \dd x^{\mu_1}\wedge\cdots .
\end{equation*}
只保留 $(\rho,\nu)$ 的反对称部分即可，因为楔积会消去对称部分：
\begin{equation*}
 \partial_\rho\partial_\nu\omega\,
 \dd x^\rho\wedge\dd x^\nu
 =\frac12(\partial_\rho\partial_\nu-\partial_\nu\partial_\rho)\omega\,
 \dd x^\rho\wedge\dd x^\nu=0.
\end{equation*}
所以 $\dd^2=0$。

再令 $\alpha^{(p)}$ 为 $p$ 阶形式。对 $\alpha^{(p)}\wedge\beta$ 的系数求导时，导数落在 $\beta$ 上所产生的新基一形式必须穿过 $p$ 个属于 $\alpha$ 的基一形式，因此得到 $(-1)^p$：
\begin{equation*}
 \dd(\alpha^{(p)}\wedge\beta)
 =(\dd\alpha^{(p)})\wedge\beta
 +(-1)^p\alpha^{(p)}\wedge\dd\beta.
\end{equation*}
两条结果合在一起即为正文\eqnrefs{eq:exterior-calculus-identities-dp68}。
\end{proof}

\begin{proposition}[曲率二形式的分量]
若 $D=\dd+A$ 且 $A=A_\mu\dd x^\mu$，则 $D^2=F\wedge$，$F=\dd A+A\wedge A=\frac12(\partial_\mu A_\nu-\partial_\nu A_\mu+[A_\mu,A_\nu])\dd x^\mu\wedge\dd x^\nu$。
\end{proposition}
\begin{proof}
对表示空间取值的零形式 $\Psi$ 连续作用两次协变微分，
\begin{align*}
 D^2\Psi
 &=\dd(\dd\Psi+\ii g_{\rm G}\mathscr C\Psi)
 +\ii g_{\rm G}\mathscr C\wedge(\dd\Psi+\ii g_{\rm G}\mathscr C\Psi).
\end{align*}
用正文\eqnrefs{eq:exterior-calculus-identities-dp68} 中的 $\dd^2=0$ 和一形式的分次 Leibniz 规则，
\begin{equation*}
 \dd(\mathscr C\Psi)=(\dd\mathscr C)\Psi-\mathscr C\wedge\dd\Psi,
\end{equation*}
于是两个含 $\mathscr C\wedge\dd\Psi$ 的交叉项相消，得到
\begin{align*}
 D^2\Psi
 &=\ii g_{\rm G}(\dd\mathscr C)\Psi
 -g_{\rm G}^2\mathscr C\wedge\mathscr C\,\Psi\\
 &=\ii g_{\rm G}
 \left(\dd\mathscr C+\ii g_{\rm G}\mathscr C\wedge\mathscr C\right)\Psi.
\end{align*}
这就给出正文\eqnrefs{eq:curvature-form-dp11}。

再分别展开 $\dd\mathscr C$ 与 $\mathscr C\wedge\mathscr C$。首先
\begin{align*}
 \dd\mathscr C
 &=(\partial_\mu\mathscr C_\nu)\dd x^\mu\wedge\dd x^\nu\\
 &=\frac12(\partial_\mu\mathscr C_\nu-\partial_\nu\mathscr C_\mu)
 \dd x^\mu\wedge\dd x^\nu.
\end{align*}
其次，矩阵系数不能交换，
\begin{align*}
 \mathscr C\wedge\mathscr C
 &=\mathscr C_\mu\mathscr C_\nu\dd x^\mu\wedge\dd x^\nu\\
 &=\frac12[\mathscr C_\mu,\mathscr C_\nu]\dd x^\mu\wedge\dd x^\nu.
\end{align*}
与 $\mathcal F=\tfrac12\mathcal F_{\mu\nu}\dd x^\mu\wedge\dd x^\nu$ 比较系数，得到正文\eqnrefs{eq:curvature-components-dp68}。
\end{proof}

\begin{proposition}[曲率的规范协变性]
固定本章 $\psi^\prime=U\psi$ 的约定，若 $D^\prime=UDU^{-1}$，则 $F^\prime=UFU^{-1}$。所有矩阵迹型不变多项式因此规范不变。
\end{proposition}
\begin{proof}
把 $\Psi'=U\Psi$ 代入正文\eqnrefs{eq:covariant-derivative-covariance-dp68}，左边为
\begin{align*}
 D'\Psi'
 &=\dd(U\Psi)+\ii g_{\rm G}\mathscr C'U\Psi\\
 &=(\dd U)\Psi+U\dd\Psi+\ii g_{\rm G}\mathscr C'U\Psi,
\end{align*}
而右边为
\begin{equation*}
 U(D\Psi)=U\dd\Psi+\ii g_{\rm G}U\mathscr C\Psi.
\end{equation*}
消去共同的 $U\dd\Psi$，并利用关系对任意 $\Psi$ 成立，
\begin{equation*}
 \dd U+\ii g_{\rm G}\mathscr C'U=\ii g_{\rm G}U\mathscr C.
\end{equation*}
右乘 $U^{-1}$ 后立即得到正文\eqnrefs{eq:finite-gauge-connection-dp11}。

由刚得到的联络变换式，协变导数满足
\[
D'=UDU^{-1}.
\]
为避免把算符乘法与普通矩阵乘法混淆，让两边作用在任意试探场 $\Psi'$ 上，并写 $\Psi'=U\Psi$。连续作用两次协变导数得到
\begin{align*}
D'^2\Psi'
&=D'\!\left(U D\Psi\right)\\
&=U D(D\Psi)\\
&=U D^2\Psi.
\end{align*}
正文对曲率的定义给出
\[
D^2\Psi=\ii g_{\rm G}\mathcal F\Psi,
\qquad
D'^2\Psi'=\ii g_{\rm G}\mathcal F'\Psi'.
\]
代入 $\Psi'=U\Psi$ 后，两种写法分别为
\begin{align*}
D'^2\Psi'
&=\ii g_{\rm G}\mathcal F' U\Psi,\\
D'^2\Psi'
&=\ii g_{\rm G}U\mathcal F\Psi.
\end{align*}
因为等式对任意 $\Psi$ 成立，且 $U$ 可逆，所以
\[
\mathcal F'U=U\mathcal F,
\qquad
\mathcal F'=U\mathcal F U^{-1},
\]
即正文\eqnrefs{eq:curvature-gauge-covariance-dp68}。
\end{proof}

\begin{theorem}[非 Abel Bianchi 恒等式]
由联络 $A$ 构造的曲率 $F=\dd A+A\wedge A$ 满足 $D_AF=\dd F+A\wedge F-F\wedge A=0$。
\end{theorem}
\begin{proof}
Bianchi 恒等式由协变导数对易子的 Jacobi 恒等式直接推出。写出
\begin{equation*}
 [D_\lambda,[D_\mu,D_\nu]]
 +[D_\mu,[D_\nu,D_\lambda]]
 +[D_\nu,[D_\lambda,D_\mu]]=0.
\end{equation*}
由正文\eqnrefs{eq:covariant-commutator-curvature-dp68}，第一项作用在任意 $\Psi$ 上为
\begin{align*}
 [D_\lambda,[D_\mu,D_\nu]]\Psi
 &=\ii g_{\rm G}[D_\lambda,\mathcal F_{\mu\nu}]\Psi\\
 &=\ii g_{\rm G}(D_\lambda\mathcal F_{\mu\nu})\Psi,
\end{align*}
其中
\begin{equation*}
 D_\lambda\mathcal F_{\mu\nu}
 =\partial_\lambda\mathcal F_{\mu\nu}
 +\ii g_{\rm G}[\mathscr C_\lambda,\mathcal F_{\mu\nu}].
\end{equation*}
对另外两项循环置换指标，Jacobi 恒等式便化为
\begin{equation*}
 \ii g_{\rm G}
 (D_\lambda\mathcal F_{\mu\nu}
 +D_\mu\mathcal F_{\nu\lambda}
 +D_\nu\mathcal F_{\lambda\mu})\Psi=0.
\end{equation*}
因为对任意 $\Psi$ 都成立，括号必须为零，即正文\eqnrefs{eq:nonabelian-bianchi-dp11}。
\end{proof}

\begin{proposition}[Wilson 线的端点协变性]
沿固定光滑路径定义的平行输运矩阵，在 $\psi^\prime=U\psi$ 的约定下满足 $W^\prime(y,x)=U(y)W(y,x)U^{-1}(x)$。闭路径输运的迹因此不变。
\end{proposition}
\begin{proof}
记沿路径的矩阵值系数为
\begin{equation*}
 A(s)=\dot x^\mu(s)\mathscr C_\mu[x(s)].
\end{equation*}
把正文\eqnrefs{eq:parallel-transport-ode-dp68} 从 $s_i$ 积到 $s$，得到 Volterra 型积分方程
\begin{equation*}
 \Psi(s)=\Psi(s_i)-\ii g_{\rm G}\int_{s_i}^{s}\dd s_1\,A(s_1)\Psi(s_1).
\end{equation*}
把右边的 $\Psi(s_1)$ 反复代回同一方程，前几阶为
\begin{align*}
 \Psi(s)=\Bigg[I
 &-\ii g_{\rm G}\int_{s_i}^{s}\dd s_1A(s_1)\\
 &+(-\ii g_{\rm G})^2
 \int_{s_i}^{s}\dd s_1\int_{s_i}^{s_1}\dd s_2\,
 A(s_1)A(s_2)\\
 &+(-\ii g_{\rm G})^3
 \int_{s_i}^{s}\dd s_1\int_{s_i}^{s_1}\dd s_2
 \int_{s_i}^{s_2}\dd s_3\,
 A(s_1)A(s_2)A(s_3)+\cdots\Bigg]\Psi(s_i).
\end{align*}
第 $n$ 阶积分域满足 $s_1\ge s_2\ge\cdots\ge s_n$；由于不同 $s$ 处的矩阵一般不交换，较大的路径参数必须排在乘积左侧。用 $\mathcal P$ 记这种顺序，级数重求和为
\begin{equation*}
 \Psi(x_f)=U_C\Psi(x_i),\qquad
 U_C=\mathcal P\exp\!\left[-\ii g_{\rm G}\int_C\mathscr C_\mu\dd x^\mu\right],
\end{equation*}
即正文\eqnrefs{eq:wilson-line-general-dp11}。

现在改变内部空间的局域基底，令 $\Psi'(x)=U(x)\Psi(x)$。同一个输运关系在新基底中满足
\begin{align*}
 \Psi'(x_f)
 &=U(x_f)\Psi(x_f)\\
 &=U(x_f)U_C\Psi(x_i)\\
 &=U(x_f)U_CU^{-1}(x_i)\Psi'(x_i).
\end{align*}
另一方面，新基底中的 Wilson 线 $U_C'$ 按正文定义满足
\begin{equation*}
 \Psi'(x_f)=U_C'\Psi'(x_i).
\end{equation*}
因为该等式对任意初始内部向量 $\Psi'(x_i)$ 都成立，比较线性算符的系数得到
\begin{equation*}
 U_C'=U(x_f)U_CU^{-1}(x_i),
\end{equation*}
即正文\eqnrefs{eq:wilson-endpoint-covariance-dp68}。
\end{proof}

\begin{proposition}[小回路与曲率]
光滑联络沿小坐标矩形的输运，相对于单位矩阵的最低非平凡项由曲率乘有向面积给出，具体符号按\eqnrefs{eq:parallel-transport-ode-dp68}与\eqnrefs{eq:small-loop-curvature-dp11}的路径次序固定；其余项至少为边长三阶。
\end{proposition}
\begin{proof}
取左下角为 $x$、两条边向量分别为 $a^\mu$ 与 $b^\nu$ 的小矩形，只保留面积阶 $ab$。四条边的短程 Wilson 线依次为
\begin{align*}
 U_1&=I-\ii g_{\rm G}\mathscr C_\mu(x)a^\mu+\order(a^2),\\
 U_2&=I-\ii g_{\rm G}[\mathscr C_\nu(x)+a^\mu\partial_\mu\mathscr C_\nu(x)]b^\nu
 +\order(b^2,a^2b),\\
 U_3&=I+\ii g_{\rm G}[\mathscr C_\mu(x)+b^\nu\partial_\nu\mathscr C_\mu(x)]a^\mu
 +\order(a^2,ab^2),\\
 U_4&=I+\ii g_{\rm G}\mathscr C_\nu(x)b^\nu+\order(b^2).
\end{align*}
闭合输运为 $U_\square=U_4U_3U_2U_1$。所有一阶项相消；Taylor 项留下
\begin{equation*}
 -\ii g_{\rm G}(\partial_\mu\mathscr C_\nu-\partial_\nu\mathscr C_\mu)a^\mu b^\nu,
\end{equation*}
矩阵乘法的二阶交叉项留下
\begin{equation*}
  +g_{\rm G}^2[\mathscr C_\mu,\mathscr C_\nu]a^\mu b^\nu.
\end{equation*}
合并并使用正文\eqnrefs{eq:curvature-components-dp68}，便得到正文\eqnrefs{eq:small-loop-curvature-dp11}。上述显示的是混合面积项；完整二阶 Taylor 展开中的纯 $a^2$、$b^2$ 项由同一路径的正反输运相消。故光滑联络下总余项为 $O(|a|^2|b|+|a||b|^2)$。
\end{proof}

\begin{proposition}[Pontryagin 形式的闭合与联络无关性]
在固定丛上，$\operatorname{tr}(F\wedge F)$ 为闭四形式，改变联络只使其增加恰当形式。在闭定向四维底空间上，按本章规范群、迹及归一化约定构造的拓扑荷不随联络连续变形变化。整数性还要求相应丛及整系数特征类归一化。
\end{proposition}
\begin{proof}
先用联络一形式 $\mathcal A$ 写曲率 $F=\dd\mathcal A+\mathcal A\wedge\mathcal A$；本章的规范势约定对应 $\mathcal A=\ii g_{\rm G}\mathscr C$，故 $F=\ii g_{\rm G}\mathcal F$。Bianchi 恒等式及迹的分次循环性给出
\[
\dd\Tr(F\wedge F)=\Tr(D_{\mathcal A}(F\wedge F))=0.
\]
对同一丛上的联络路径 $\mathcal A_t$，差及其导数为全局伴随丛值一形式。由 $\dot F_t=D_{\mathcal A_t}\dot{\mathcal A}_t$ 得
\begin{align*}
\frac{\dd}{\dd t}\Tr(F_t\wedge F_t)
&=2\Tr(D_{\mathcal A_t}\dot{\mathcal A}_t\wedge F_t)\\
&=2\dd\Tr(\dot{\mathcal A}_t\wedge F_t).
\end{align*}
积分得两种联络的四形式之差为恰当形式；闭定向底空间上的积分由 Stokes 定理不变。以下再核对分量归一化；整数性仍来自整体特征类条件。写
\begin{equation*}
 \mathcal F=\frac12\mathcal F_{\mu\nu}^aT^a
 \dd x^\mu\wedge\dd x^\nu,
 \qquad
 \Tr(T^aT^b)=\frac12\delta^{ab}.
\end{equation*}
于是
\begin{align*}
 \Tr(\mathcal F\wedge\mathcal F)
 &=\frac18\mathcal F_{\mu\nu}^a\mathcal F_{\rho\sigma}^a
 \epsilon^{\mu\nu\rho\sigma}\dd^4 x.
\end{align*}
再用对偶场强定义
\begin{equation*}
 {}^\star\!\mathcal F^{a\mu\nu}
 =\frac12\epsilon^{\mu\nu\rho\sigma}\mathcal F_{\rho\sigma}^a,
\end{equation*}
可写成
\begin{equation*}
 \Tr(\mathcal F\wedge\mathcal F)
 =\frac14\mathcal F_{\mu\nu}^a
 {}^\star\!\mathcal F^{a\mu\nu}\dd^4 x.
\end{equation*}
代回正文\eqnrefs{eq:pontryagin-form-definition-dp68}，得到\eqnrefs{eq:pontryagin-number-dp11} 中的分量归一化 $g_{\rm G}^2/(32\pi^2)$。其中 $Q_{\rm top}\in\mathbb Z$ 的部分依赖正文已经说明的边界条件与整体拓扑定理，不属于这个代数推导的结论。
\end{proof}

\begin{proposition}[矩阵群的 Lie 括号]
矩阵 Lie 群的切空间在矩阵交换子下封闭。对 $X,Y\in\mathfrak g$，$e^{tX}e^{sY}e^{-tX}e^{-sY}=I+ts[X,Y]+O(|t|^2|s|+|t||s|^2)$。
\end{proposition}
\begin{proof}
先沿固定切向方向 $X=\dot g(0)$ 考察单参数子群。由群关系 $g(t+s)=g(t)g(s)$ 对 $s$ 求导并令 $s=0$，得到
\begin{align*}
 \dot g(t)
 &=g(t)\left.\frac{\dd g(s)}{\dd s}\right|_{s=0}\\
 &=g(t)X.
\end{align*}
交换 $t,s$ 的角色，同理由 $g(t+s)=g(s)g(t)$ 得 $\dot g(t)=Xg(t)$，所以 $X$ 与该单参数子群上的 $g(t)$ 对易。设
\begin{equation*}
 g(t)=\sum_{n=0}^{\infty}t^nG_n,
\end{equation*}
代入 $\dot g=Xg$ 并比较 $t^n$ 的系数，得到
\begin{equation*}
 (n+1)G_{n+1}=XG_n,\qquad G_0=I.
\end{equation*}
递推给出 $G_n=X^n/n!$，于是
\begin{equation*}
 g(t)=I+tX+\frac{t^2X^2}{2!}+\cdots=\ee^{tX},
\end{equation*}
即正文\eqnrefs{eq:one-param-exp-dp11}。

接着取另一个切向量 $Y\in\mathfrak g=T_eG$，并选一条群曲线 $h(t)$ 满足 $h(0)=I$、$\dot h(0)=Y$。由刚得到的 $\ee^{sX}\in G$ 以及群乘法和取逆的封闭性，
\begin{equation*}
 h_s(t)=\ee^{sX}h(t)\ee^{-sX}\in G.
\end{equation*}
因此它在 $t=0$ 的切向量仍属于单位元切空间：
\begin{equation*}
 Y(s)=\left.\frac{\dd h_s(t)}{\dd t}\right|_{t=0}
 =\ee^{sX}Ye^{-sX}\in\mathfrak g.
\end{equation*}
对 $s$ 再求导，
\begin{align*}
 \frac{\dd Y(s)}{\dd s}
 &=Xe^{sX}Ye^{-sX}-\ee^{sX}Ye^{-sX}X,\\
 \left.\frac{\dd Y(s)}{\dd s}\right|_{s=0}
 &=XY-YX.
\end{align*}
因为 $Y(s)$ 是向量空间 $\mathfrak g$ 内的光滑曲线，其导数仍属于 $\mathfrak g$，所以
\begin{equation*}
 XY-YX\in\mathfrak g,
\end{equation*}
从而得到正文\eqnrefs{eq:matrix-lie-bracket-math-dp76}。
\end{proof}

\begin{proposition}[表示的微分保持括号]
若 $\rho:G\to GL(V)$ 为光滑群表示，则其微分满足 $\rho_*([X,Y])=[\rho_*(X),\rho_*(Y)]$，从而给出 Lie 代数表示。
\end{proposition}
\begin{proof}
需要验证的关键点是：群表示在单位元处求导以后，仍然保持 Lie 括号。取两个抽象生成方向 $X,Y\in\mathfrak g$。群中的共轭变换满足
\[
\operatorname{Ad}_{\ee^{sX}}Y=\ee^{sX}Ye^{-sX}.
\]
表示保持群乘法，因此
\begin{align*}
D_R(\ee^{sX})D_R(\ee^{tY})D_R(\ee^{-sX})
&=D_R\!\left(\ee^{sX}\ee^{tY}\ee^{-sX}\right).
\end{align*}
先对 $t$ 在 $t=0$ 求导，得到
\begin{align*}
D_R(\ee^{sX})\,dD_R(Y)\,D_R(\ee^{-sX})
&=dD_R\!\left(\operatorname{Ad}_{\ee^{sX}}Y\right).
\end{align*}
再对 $s$ 在 $s=0$ 求导。左边给矩阵对易子，右边给抽象伴随作用的导数：
\begin{align*}
[dD_R(X),dD_R(Y)]
&=dD_R([X,Y]_{\mathfrak g}).
\end{align*}
现在令 $X=X_a$、$Y=X_b$，并代入正文\eqnrefs{eq:lie-algebra-general-dp11}：
\begin{align*}
[dD_R(X_a),dD_R(X_b)]
&=f_{ab}{}^c\,dD_R(X_c).
\end{align*}
对酉表示，正文已定义 $dD_R(X_a)=-\ii T_R^a$，于是
\begin{align*}
[-\ii T_R^a,-\ii T_R^b]
&=f_{ab}{}^c(-\ii T_R^c),\\
-[T_R^a,T_R^b]
&=-\ii f_{ab}{}^cT_R^c.
\end{align*}
两边乘以 $-1$，得到正文\eqnrefs{eq:rep-lie-algebra-dp11}：
\[
[T_R^a,T_R^b]=\ii f_{ab}{}^cT_R^c.
\]
\end{proof}

\begin{lemma}[基本表示的完整性恒等式]
设 $SU(N)$ 的 Hermitian 无迹生成元满足 $\operatorname{tr}(T^aT^b)=T_F\delta^{ab}$。则 $\sum_a(T^a)_{ij}(T^a)_{kl}=T_F(\delta_{il}\delta_{jk}-N^{-1}\delta_{ij}\delta_{kl})$。
\end{lemma}
\begin{proof}
正文已经固定基本表示生成元的无迹性和迹归一化。四指标生成元和只能由两个 Kronecker 张量构成；用两种收缩固定它们的系数，就得到正文\eqnrefs{eq:suN-fierz-math-dp68}。设 $SU(N)$ 生成元满足
\begin{equation*}
 \Tr(T^aT^b)=T_F\delta^{ab},
 \qquad \Tr T^a=0.
\end{equation*}
把所有无迹厄米生成元与单位矩阵一起看成 $N\times N$ 矩阵空间的一组正交基。定义
\begin{equation*}
 X_{ij;kl}\equiv\sum_{a=1}^{N^2-1}(T^a)_{ij}(T^a)_{kl}.
\end{equation*}
由于 $i,l$ 与 $j,k$ 分别形成基本表示--反基本表示指标对，$SU(N)$ 不变量只能由
\begin{equation*}
 \delta_{il}\delta_{jk},
 \qquad
 \delta_{ij}\delta_{kl}
\end{equation*}
组成，所以写
\begin{equation*}
 X_{ij;kl}=A\delta_{il}\delta_{jk}+B\delta_{ij}\delta_{kl}.
\end{equation*}

第一组收缩取 $i=j$ 并求和：
\begin{align*}
 \sum_iX_{ii;kl}
 &=\sum_a\Tr(T^a)(T^a)_{kl}=0,\\
 &=A\sum_i\delta_{il}\delta_{ik}
 +B\sum_i\delta_{ii}\delta_{kl},\\
 &=(A+NB)\delta_{kl}.
\end{align*}
所以
\begin{equation*}
 B=-\frac AN.
\end{equation*}
第二组收缩取 $i=l$、$j=k$ 并对 $i,j$ 求和：
\begin{align*}
 \sum_{ij}X_{ij;ji}
 &=\sum_a\Tr(T^aT^a),\\
 &=T_F(N^2-1).
\end{align*}
而拟设右边给
\begin{align*}
 \sum_{ij}
 \left(A\delta_{ii}\delta_{jj}+B\delta_{ij}\delta_{ji}\right)
 &=AN^2+BN,\\
 &=A(N^2-1).
\end{align*}
故 $A=T_F$，从而
\begin{equation*}
 {
 \sum_a(T^a)_{ij}(T^a)_{kl}
 =T_F\left(
 \delta_{il}\delta_{jk}
 -\frac1N\delta_{ij}\delta_{kl}
 \right).}
\end{equation*}

检查：令 $k=j$ 并求和 $j$，左边是 $(\sum_aT^aT^a)_{il}=C_F\delta_{il}$；右边给
\begin{equation*}
 T_F\left(N-\frac1N\right)\delta_{il}
 =T_F\frac{N^2-1}{N}\delta_{il},
\end{equation*}
所以 $C_F=T_F(N^2-1)/N$，对 $T_F=1/2$ 恢复 $C_F=(N^2-1)/(2N)$。
\end{proof}

\begin{corollary}[基本与伴随 Casimir]
采用 $T_F=1/2$ 的生成元归一化，$SU(N)$ 基本表示满足 $\sum_aT^aT^a=C_FI$，$C_F=(N^2-1)/(2N)$；伴随表示满足 $f^{acd}f^{bcd}=C_A\delta^{ab}$，$C_A=N$。
\end{corollary}
\begin{proof}
伴随 Casimir 可以不写任何显式 $SU(N)$ 矩阵，而由刚得到的完备关系求出。比较同一个对易子迹的两种计算即可固定 $C_A$。从基本对易子
\begin{equation*}
 [T^a,T^b]=\ii f^{abc}T^c
\end{equation*}
出发。考虑
\begin{equation*}
 \Tr([T^a,T^c][T^b,T^c]),
\end{equation*}
并对重复的 $c$ 求和。一方面直接代入结构常数：
\begin{align*}
 \sum_c\Tr([T^a,T^c][T^b,T^c])
 &=\sum_c\Tr(\ii f^{acd}T^d\,\ii f^{bce}T^e),\\
 &=-T_F f^{acd}f^{bcd},\\
 &=-T_F C_A\delta^{ab}.
\end{align*}
另一方面展开对易子：
\begin{align*}
 \sum_c\Tr([T^a,T^c][T^b,T^c])
 ={}&\sum_c\operatorname{Tr}\!\left[
 T^aT^cT^bT^c-T^aT^cT^cT^b\right.\\
 &\left.\hspace{2.5cm}-T^cT^aT^bT^c+T^cT^aT^cT^b
 \right].
\end{align*}
逐项用 Fierz/完备关系。对任意矩阵 $X$，由 Fierz 恒等式可得
\begin{equation*}
 \sum_c T^cXT^c
 =T_F\left[(\Tr X)I-\frac1N X\right].
\end{equation*}
因此
\begin{align*}
 \sum_c\Tr(T^aT^cT^bT^c)
 &=-\frac{T_F}{N}\Tr(T^aT^b),\\
 \sum_c\Tr(T^cT^aT^cT^b)
 &=-\frac{T_F}{N}\Tr(T^aT^b).
\end{align*}
另外 $\sum_cT^cT^c=C_FI$，故中间两项各为
\begin{equation*}
 -C_F\Tr(T^aT^b).
\end{equation*}
合并：
\begin{align*}
 \sum_c\Tr([T^a,T^c][T^b,T^c])
 &=-2\left(C_F+\frac{T_F}{N}\right)
 \Tr(T^aT^b),\\
 &=-2T_F\left(C_F+\frac{T_F}{N}\right)\delta^{ab}.
\end{align*}
取 $T_F=1/2$、$C_F=(N^2-1)/(2N)$：
\begin{equation*}
 2\left(C_F+\frac{T_F}{N}\right)
 =2\left(\frac{N^2-1}{2N}+\frac1{2N}\right)=N.
\end{equation*}
与第一种计算比较即得
\begin{equation*}
 C_A=N.
\end{equation*}
这条推导清楚展示了伴随 Casimir 与基本完备关系之间的关系。
\end{proof}

\begin{theorem}[Sokhotski--Plemelj 分布公式]
在实直线上的缓增分布意义下，$\lim_{\epsilon\downarrow0}(x\pm\ii\epsilon)^{-1}=\operatorname{PV}(1/x)\mp\ii\pi\delta(x)$。
\end{theorem}
\begin{proof}
正文\eqnrefs{eq:plemelj-split-r13} 已把复分母分成实部和虚部两个普通函数。$\epsilon\to0^+$ 时，实部按正文\eqnrefs{eq:principal-value-definition-dp68} 趋于 Cauchy 主值，虚部则形成单位面积、宽度趋零的 Lorentz 峰；分别取这两个分布极限即可得到正文\eqnrefs{eq:sokhotski-plemelj-dp68}。

先令
\begin{equation*}
 \delta_\epsilon(x)
 =\frac1\pi\frac{\epsilon}{x^2+\epsilon^2}.
\end{equation*}
变量代换 $x=\epsilon u$ 给
\begin{equation*}
 \int\dd x\,\delta_\epsilon(x)=1,
\end{equation*}
且对任意光滑测试函数 $f$，
\begin{equation*}
 \int\dd x\,\delta_\epsilon(x)f(x)
 =\frac1\pi\int\frac{\dd u}{1+u^2}f(\epsilon u)
 \longrightarrow f(0).
\end{equation*}
所以 $\delta_\epsilon\to\delta$。另一方面，把
$x/(x^2+\epsilon^2)$ 的正负半轴配对后，极限正是正文\eqnrefs{eq:principal-value-definition-dp68}。代回正文\eqnrefs{eq:plemelj-split-r13} 即得正文\eqnrefs{eq:sokhotski-plemelj-dp68}。
\end{proof}

\begin{proposition}[指数核围道的方向与留数]
对频率指数 $e^{-\ii\omega t}$，$t>0$ 时应在下半平面闭合，闭合围道顺时针给 $-2\pi\ii$ 倍留数和；$t<0$ 时在上半平面闭合给 $+2\pi\ii$ 倍留数和。弧积分消失要求被积函数满足对应的 Jordan 衰减条件。
\end{proposition}
\begin{proof}
正文\eqnrefs{eq:jordan-exponential-r13} 已告诉我们指数在复平面的哪一侧衰减。还需要再检查大半圆积分确实消失以及闭合曲线的取向：$\tau>0$ 时下半圆顺时针，$\tau<0$ 时上半圆逆时针。两点合起来便得到正文\eqnrefs{eq:contour-orientation-residue-dp68} 的正负号。

设大 $|z|$ 时 $|F(z)|\le C/|z|^{1+\delta}$，$\delta>0$。对 $\tau>0$ 取下半圆 $z=Re^{-\ii\theta}$，则
\begin{equation*}
 |\ee^{-\ii z\tau/\hbar}|=\ee^{-R\tau\sin\theta/\hbar}\le1,
\end{equation*}
从而
\begin{equation*}
 \left|\int_{C_R^-}\dd z\,F(z)\ee^{-\ii z\tau/\hbar}\right|
 \le\frac{\pi C}{R^\delta}\to0.
\end{equation*}
下半圆从 $+R$ 绕回 $-R$，闭合曲线顺时针，因此留数定理给实轴积分
$-2\pi\ii$ 乘下半平面留数；除以正文 Fourier 约定 中的 $2\pi$ 即得\eqnrefs{eq:contour-orientation-residue-dp68} 第一行。$\tau<0$ 时同理向上闭合，曲线逆时针，得到第二行。
\end{proof}
